
The \hyperref[module-wext]{\texttt{wext}} module manages AIMS build
extensions such as MPI or Scalapack.

To make a new build extension, create a file called
\texttt{wext\_\textless{}name\textgreater{}.py} in the \texttt{src}
directory where \texttt{\textless{}name\textgreater{}} is the name of
your extension. This Python file can be empty, in which case the
extension has no effect besides adding an option to the Waf options (see
\texttt{./waf\ -\/-help}). The extension is then activated by calling
Waf with this option. Functionality is added by defining variables and
functions (all optional) in
\texttt{wext\_\textless{}name\textgreater{}.py}:

\begin{description}
\item[help]
A string with a help for the build extension which appears in
\texttt{./waf\ -\/-help}.
\item[require]
A list of build extensions that the extension requires if activated.
\item[stub(ctx)]
Function that takes a \hyperref[wext.Context]{\texttt{Context}} (see
below) as an argument and applies actions (such as adding source files)
that should happen if the extension is not active.
\item[extension(ctx)]
Takes a \hyperref[wext.Context]{\texttt{Context}} as argument and
applies actions that should happen when the extension is active.
\end{description}

For example, the Scalapack extension is defined as:

\begin{verbatim}
help = 'compile aims with Scalapack'
require = ['mpi']

def stub(ctx):
    ctx.add_files('fc', 'scalapack_stub.f90')

def extension(ctx):
    ctx.add_link_flag(ctx.env.SCALAPACK)
\end{verbatim}

In common language, Scalapack requires MPI to work. The Scalapack stub
adds an appropriate stub file (compiled with Fortran compiler -- `fc').
And finally if Scalapack is activated, a link flag is added from the
\texttt{SCALAPACK} environment variable. For full list of avaialble
actions, see documentation of \hyperref[wext.Context]{\texttt{Context}}.

\begin{description}
\item[ \emph{class} \texttt{wext.}\texttt{Context}{(}\emph{wctx}{)}]
Provides interface to modify a build process.

\begin{description}
\item[ \texttt{add\_feature}{(}\emph{feat}, \emph{files}{)}]
Add a feature to files.
\end{description}

\begin{description}
\item[ \texttt{add\_files}{(}\emph{feat}, \emph{files}{)}]
Add files with features to the build process.
\end{description}

\begin{description}
\item[ \texttt{add\_link\_flag}{(}\emph{flag}{)}]
Add a link flag.
\end{description}

\begin{description}
\item[ \texttt{change\_ext}{(}\emph{fext}{)}]
Change file name extension (including the dot) of the resulting binary.
\end{description}

\begin{description}
\item[ \texttt{env} \emph{= None}]
Provides access to the variable environment.
\end{description}

\begin{description}
\item[ \texttt{replace\_files}{(}\emph{newfiles}, \emph{oldfiles},
\emph{feat=None}{)}]
Replace old files with new files.

Optionally add a feature. All old files have to have the same features.
\end{description}

\begin{description}
\item[ \texttt{wctx} \emph{= None}]
Provides access to the Waf build configuration context.
\end{description}
\end{description}

